\documentclass{article}
\usepackage{amsmath, amssymb, amsthm}
\usepackage{hyperref}
\usepackage{geometry}
\geometry{margin=1in}

\title{Erd\H{o}s Problem \#629}
\date{Last updated: 2025-08-31}
\author{Source: erdosproblems.com}

\begin{document}
\maketitle

\noindent\textbf{Status:} OPEN \quad \textbf{No prize}

\noindent\textbf{Tags:} graph theory, chromatic number

\medskip
\noindent\textbf{Problem Statement:}

The list chromatic number $\chi_L(G)$ is defined to be the minimal $k$ such that for any assignment of a list of $k$ colours to each vertex of $G$ (perhaps different lists for different vertices) a colouring of each vertex by a colour on its list can be chosen such that adjacent vertices receive distinct colours.

Determine the minimal number of vertices $n(k)$ of a bipartite graph $G$ such that $\chi_L(G)>k$.

\bigskip
\noindent\textbf{Remarks:}

A problem of Erd\H{o}s, Rubin, and Taylor \cite{ERT80}, who proved that\[2^{k-1}<n(k) <k^22^{k+2}.\]They also prove that if $m(k)$ is the size of the smallest family of $k$-sets without property B (i.e. the smallest number of $k$-sets in a graph with chromatic number $3$) then $m(k)\leq n(k)\leq m(k+1)$. Bounds on $m(k)$ are the subject of [901]. The lower bounds on $m(k)$ by Radhakrishnan and Srinivasan \cite{RaSr00} imply that\[2^k \left(\frac{k}{\log k}\right)^{1/2}\ll n(k).\]Erd\H{o}s, Rubin, and Taylor \cite{ERT80} proved $n(2)=6$ and Hanson, MacGillivray, and Toft \cite{HMT96} proved $n(3)=14$ and\[n(k) \leq kn(k-2)+2^k.\]See also the entry in the graphs problem collection.

\bigskip
\noindent\textbf{References:}

[ERT80] Erd\H{o}s, Paul and Rubin, Arthur L. and Taylor, Herbert, Choosability in graphs. (1980), 125-157.
 
 [HMT96] Hanson, Denis and MacGillivray, Gary and Toft, Bjarne, Choosability of bipartite graphs. Ars Combin. (1996), 183-192.
 
 [RaSr00] Radhakrishnan, Jaikumar and Srinivasan, Aravind, Improved bounds and algorithms for hypergraph {$2$}-coloring. Random Structures Algorithms (2000), 4--32.

\bigskip
\noindent\small{Source: \url{https://www.erdosproblems.com/629}}
\end{document}
